\documentclass[a4paper]{article}

%------------------------------------------------
% Packages
%------------------------------------------------

\usepackage{enumitem}
\usepackage[margin=2.5cm]{geometry}
\usepackage{amsmath,amsfonts,amssymb}
\usepackage[charter]{mathdesign}
\usepackage{graphicx}
\usepackage{array}
\usepackage{fancyhdr}
\usepackage{polyglossia}

%------------------------------------------------
% Langues et Polices
%------------------------------------------------

\setdefaultlanguage[
    calendar=gregorian,
    numerals=maghrib
]{arabic}

\setotherlanguage{english}

% إذا لم تكن الخطوط مثبتة في الجهاز، تُؤخذ من المجلد ../fonts/
\IfFontExistsTF{Amiri}{%
    \newfontfamily\arabicfont[Script=Arabic, Scale=1.15]{Amiri}
}{%
    \newfontfamily\arabicfont[
        Script=Arabic, Scale=1.15,
        Path=../fonts/, Extension=.ttf,
        UprightFont=*-Regular, BoldFont=*-Bold,
        ItalicFont=*-Italic, BoldItalicFont=*-BoldItalic
    ]{Amiri}
}

\IfFontExistsTF{AlBattar}{%
    \newfontfamily\arabicfontsf[Script=Arabic, Scale=1.2]{AlBattar}
    % خط العناوين (يعمل حتى بعد استعمال \textenglish)
    \newfontfamily\khatt[Script=Arabic, Scale=1.2]{AlBattar}
}{%
    \newfontfamily\arabicfontsf[
        Script=Arabic, Scale=1.2, Path=../fonts/,
        BoldFont=ae_AlBattar.ttf
    ]{ae_AlBattar.ttf}
    % خط العناوين (يعمل حتى بعد استعمال \textenglish)
    \newfontfamily\khatt[
        Script=Arabic, Scale=1.2, Path=../fonts/,
        BoldFont=ae_AlBattar.ttf
    ]{ae_AlBattar.ttf}
}

%------------------------------------------------
% Configuration de la page
%------------------------------------------------

\renewcommand{\footrulewidth}{1pt}
\renewcommand{\headrulewidth}{1pt}

\rhead{\khatt المدرسة العليا للأساتذة ببوزريعة}
\lhead{\khatt ملحقة بومرداس}

\cfoot{}
\lfoot{\khatt الأستاذ: د. بن عبد الرزاق يوسف}

\parindent 0pt
\setlength{\headsep}{15pt}
\setlength{\parskip}{3pt}

\pagestyle{fancy}

\setlist[enumerate,1]{
    label=\arabic*),
    itemsep=5pt,
    topsep=3pt
}

\setlist[itemize]{
    itemsep=3pt,
    topsep=3pt
}

\arraycolsep=3pt

%================================================
% Début du document
%================================================

\begin{document}

%================================================
% TITRE
%================================================

\begin{center}

{\Large\khatt\textbf{سلسلة تمارين رقم 7}}

\vspace{0.15cm}

{\Large\khatt\textbf{
المخبر: بناء مساعد ذكاء اصطناعي وسير عمل
}}

\vspace{0.10cm}

{\khatt\textbf{
إنشاء ــ تهيئة ــ معرفة ــ اختبار ــ تحسين ــ سير عمل
}}

\end{center}

\rule{\textwidth}{1pt}

\begin{center}

{\khatt\textbf{
الفصل 7: تصميم مساعد ذكاء اصطناعي وسير عمل
}}

\end{center}

\rule{\textwidth}{1pt}


%================================================
% الوضعية المحفزة
%================================================

\section*{\khatt\underline{الوضعية المحفزة}}

يستعمل الأستاذ الذكاء الاصطناعي بشكل متكرر في
تلخيص الدروس، وشرح المفاهيم، وإنشاء الأسئلة
التدريبية.

لكن كتابة نفس التعليمات وإضافة نفس الملفات في كل
محادثة قد تكون عملية متعبة.

لذلك قرر الأستاذ إنشاء مساعد ذكاء اصطناعي مخصص
لمادة دراسية معينة، ثم تنظيم المهام الكبيرة في
عدة مراحل.

\textbf{السؤال:}

كيف يمكن بناء مساعد ذكاء اصطناعي مخصص، تزويده
بالمعلومات اللازمة، اختباره وتحسينه، ثم بناء
سير عمل يقوم بمهمة كبيرة على عدة مراحل؟


%================================================
% أهداف المخبر
%================================================

\section*{\khatt\underline{ما سننجزه في هذا المخبر}}

في هذا المخبر سنبني مساعدًا تعليميًا خطوة بخطوة.

في نهاية العمل سيكون لدينا:

\begin{itemize}

    \item مساعد ذكاء اصطناعي مخصص.

    \item تعليمات واضحة للمساعد.

    \item ملف مرجعي يستعمله المساعد.

    \item مجموعة من الاختبارات.

    \item تعليمات محسنة بعد الاختبار.

    \item سير عمل مكون من عدة مراحل.

\end{itemize}

سنطبق الفكرة باستعمال
\textenglish{ChatGPT} أو \textenglish{Claude}.


%================================================
% النشاط الأول
%================================================

\section*{\khatt\underline{النشاط الأول: تحديد المساعد الذي سنبنيه}}

في هذا المخبر سنبني مساعدًا باسم:

\begin{center}

\fbox{
\parbox{10cm}{
\centering
\textbf{مساعد مادة الذكاء الاصطناعي}
}
}

\end{center}

وظيفة المساعد هي مساعدة الطالب على فهم الدروس
باستعمال لغة بسيطة وأمثلة مناسبة.

\subsection*{\khatt الخطوة 1: تحديد المهمة}

المساعد الذي سنبنيه يستطيع القيام بالمهام التالية:

\begin{itemize}

    \item شرح مفهوم للطالب.

    \item تلخيص جزء من الدرس.

    \item إنشاء أسئلة تدريبية.

\end{itemize}

\textbf{المهمة الرئيسية للمساعد:}

\[
\boxed{\hspace{15cm}}
\]

\subsection*{\khatt الخطوة 2: تحديد المستخدم}

المستخدم المستهدف هو:

\[
\boxed{\hspace{15cm}}
\]

\subsection*{\khatt الخطوة 3: تحديد المدخلات}

ما الذي يمكن للطالب إدخاله إلى المساعد؟

\[
\boxed{\hspace{15cm}}
\]

مثال:

\begin{center}

سؤال ــ نص درس ــ مفهوم ــ فقرة

\end{center}

\subsection*{\khatt الخطوة 4: تحديد المخرجات}

ماذا يجب أن ينتج المساعد؟

\[
\boxed{\hspace{15cm}}
\]


%================================================
% النشاط الثاني
%================================================

\section*{\khatt\underline{النشاط الثاني (واجب منزلي): تحضير الملف المرجعي}}

\textbf{يُنجز هذا النشاط في المنزل ويُسلَّم في الحصة القادمة.}

لكي يستطيع المساعد الاستفادة من محتوى المادة،
سنضيف إليه ملفًا مرجعيًا.

\subsection*{\khatt الخطوة 1: اختيار الملف}

اختر ملفًا قصيرًا متعلقًا بالمادة.

يمكن أن يكون:

\begin{itemize}

    \item ملف PDF.

    \item ملف Word.

    \item ملف نصي.

    \item محاضرة.

    \item ملخص.

\end{itemize}

\textbf{مثال:}

\begin{center}

\textenglish{Introduction to Artificial Intelligence.pdf}

\end{center}

\subsection*{\khatt الخطوة 2: فحص الملف}

قبل رفع الملف، افتحه وتحقق من:

\begin{itemize}

    \item وجود عنوان واضح.

    \item وجود معلومات مرتبطة بالمادة.

    \item عدم كون الملف فارغًا.

    \item إمكانية قراءة النص.

\end{itemize}

\textbf{اسم الملف الذي اخترته:}

\[
\boxed{\hspace{15cm}}
\]

\textbf{موضوع الملف:}

\[
\boxed{\hspace{15cm}}
\]


%================================================
% النشاط الثالث
%================================================

\section*{\khatt\underline{النشاط الثالث: كتابة تعليمات المساعد}}

قبل إنشاء المساعد، سنكتب التعليمات التي
ستحدد طريقة عمله.

التعليمات التالية هي التعليمات الأساسية التي
سنستعملها في المخبر.

\begin{center}

\fbox{
\parbox{13cm}{

\textbf{الدور:}

أنت مساعد تعليمي متخصص في مادة الذكاء الاصطناعي.

\vspace{0.15cm}

\textbf{المهمة:}

ساعد الطلاب على فهم الدروس بطريقة بسيطة
ومنظمة.

\vspace{0.15cm}

\textbf{طريقة العمل:}

اشرح المفاهيم خطوة بخطوة، واستعمل أمثلة بسيطة
عندما يكون ذلك مناسبًا.

\vspace{0.15cm}

\textbf{المراجع:}

اعتمد على الملفات التي يضيفها الأستاذ عندما
يكون السؤال متعلقًا بالمادة.

\vspace{0.15cm}

\textbf{القيود:}

لا تخترع معلومات غير موجودة في المراجع.

إذا لم تجد المعلومة المطلوبة، صرّح بذلك بوضوح.

\vspace{0.15cm}

\textbf{شكل الإجابة:}

استعمل عناوين واضحة ونقاطًا وأمثلة بسيطة.

\vspace{0.15cm}

\textbf{التعامل مع السؤال غير الواضح:}

اطلب من الطالب توضيح السؤال قبل إعطاء إجابة
غير مؤكدة.

}
}

\end{center}

\textbf{ملاحظة:}

هذه التعليمات سننسخها لاحقًا ونضعها داخل إعدادات
المساعد.


%================================================
% النشاط الرابع
%================================================

\section*{\khatt\underline{النشاط الرابع: إنشاء المساعد في ChatGPT}}

\textbf{الهدف:}

إنشاء مساعد مخصص باستعمال إعدادات
\textenglish{GPT} إذا كانت الخاصية متاحة في
الحساب أو مساحة العمل.

\subsection*{\khatt الخطوة 1: فتح ChatGPT}

\begin{enumerate}

    \item افتح متصفح الإنترنت.

    \item افتح \textenglish{ChatGPT}.

    \item سجل الدخول إلى حسابك.

\end{enumerate}

بعد تسجيل الدخول يجب أن تظهر لك واجهة ChatGPT.


%------------------------------------------------

\subsection*{\khatt الخطوة 2: فتح GPTs}

ابحث في واجهة ChatGPT عن:

\begin{center}

\textenglish{GPTs}

\end{center}

ثم افتح القسم.

ابحث عن الزر:

\begin{center}

\textenglish{Create}
\end{center}

واضغط عليه.

\textbf{إذا لم يظهر الخيار:}

انتقل إلى جزء Claude في هذا المخبر.


%------------------------------------------------

\subsection*{\khatt الخطوة 3: فتح إعدادات المساعد}

بعد الضغط على
\textenglish{Create}
ستظهر واجهة إنشاء المساعد.

ابحث عن إعدادات التكوين، مثل:

\begin{center}

\textenglish{Configure}
\end{center}

سنقوم الآن بتعبئة الإعدادات واحدة تلو الأخرى.


%------------------------------------------------

\subsection*{\khatt الخطوة 4: كتابة اسم المساعد}

في خانة الاسم اكتب:

\begin{center}

\textbf{مساعد مادة الذكاء الاصطناعي}

\end{center}


%------------------------------------------------

\subsection*{\khatt الخطوة 5: كتابة وصف المساعد}

في خانة الوصف اكتب:

\begin{center}

\fbox{
\parbox{11cm}{
\centering
مساعد تعليمي يساعد الطلاب على فهم دروس الذكاء
الاصطناعي وشرح المفاهيم وإنشاء أسئلة تدريبية.
}
}

\end{center}


%------------------------------------------------

\subsection*{\khatt الخطوة 6: إضافة التعليمات}

ابحث عن الحقل:

\begin{center}

\textenglish{Instructions}

\end{center}

هذا الحقل يحدد طريقة عمل المساعد.

انسخ التعليمات التالية والصقها كاملة:

\begin{center}

\fbox{
\parbox{13cm}{

أنت مساعد تعليمي متخصص في مادة الذكاء الاصطناعي.

مهمتك مساعدة الطلاب على فهم الدروس بطريقة بسيطة
ومنظمة.

اشرح المفاهيم خطوة بخطوة واستعمل أمثلة بسيطة.

اعتمد على الملفات التي يضيفها الأستاذ عندما يكون
السؤال متعلقًا بالمادة.

لا تخترع معلومات غير موجودة في المراجع.

إذا لم تجد المعلومة المطلوبة، صرّح بذلك بوضوح.

نظم الإجابات باستعمال العناوين والنقاط والأمثلة.

عندما يكون السؤال غير واضح، اطلب توضيحًا من الطالب.

}
}

\end{center}


%------------------------------------------------

\subsection*{\khatt الخطوة 7: إضافة رسالة بداية}

ابحث عن إعداد:

\begin{center}

\textenglish{Conversation Starters}

\end{center}

أضف مثلًا:

\begin{enumerate}

    \item اشرح لي مفهوم الذكاء الاصطناعي.

    \item لخص لي هذا الدرس.

    \item أنشئ لي أسئلة تدريبية.

\end{enumerate}


%------------------------------------------------

\subsection*{\khatt الخطوة 8: إضافة الملف}

ابحث عن:

\begin{center}

\textenglish{Knowledge}

\end{center}

ثم أضف الملف الذي اخترته في النشاط الثاني.

انتظر حتى ينتهي رفع الملف.

\textbf{تحقق:}

يجب أن يظهر اسم الملف ضمن الملفات المضافة.


%------------------------------------------------

\subsection*{\khatt الخطوة 9: الأدوات}

إذا كانت هناك أدوات متاحة، لا تقم بتفعيلها
عشوائيًا.

اسأل أولًا:

\begin{center}

هل يحتاج هذا المساعد فعلًا إلى هذه الأداة؟

\end{center}

مثلاً:

\begin{center}

$\square$ البحث في الويب

\hspace{0.8cm}

$\square$ تحليل البيانات

\hspace{0.8cm}

$\square$ إنشاء الصور

\end{center}

\begin{center}

$\square$ لا توجد أداة إضافية

\end{center}

\textbf{سبب الاختيار:}

\[
\boxed{\hspace{15cm}}
\]


%================================================
% النشاط الخامس
%================================================

\section*{\khatt\underline{النشاط الخامس: اختبار المساعد في ChatGPT}}

بعد إنشاء المساعد، لا نفترض أنه يعمل بشكل
صحيح.

يجب اختباره.

\subsection*{\khatt الاختبار 1: سؤال بسيط}

افتح منطقة الاختبار أو المحادثة.

اكتب:

\begin{center}

\fbox{
\parbox{11cm}{
ما هو الذكاء الاصطناعي؟

اشرحه لطالب مبتدئ.
}
}

\end{center}

اضغط على إرسال.

راقب النتيجة.

هل استعمل المساعد لغة بسيطة؟

\begin{center}

$\square$ نعم
\qquad
$\square$ لا

\end{center}


%------------------------------------------------

\subsection*{\khatt الاختبار 2: استعمال الملف}

اكتب:

\begin{center}

\fbox{
\parbox{11cm}{
اعتمد على الملف المرفق.

اذكر أهم ثلاثة مفاهيم في الدرس،
واشرح كل مفهوم في سطرين.
}
}

\end{center}

قارن الإجابة بالمعلومات الموجودة في الملف.


%------------------------------------------------

\subsection*{\khatt الاختبار 3: سؤال بمعلومات ناقصة}

اكتب:

\begin{center}

\fbox{
\parbox{11cm}{
اشرح لي هذا الموضوع بالتفصيل.
}
}

\end{center}

لاحظ هل سيطلب المساعد معلومات إضافية.


%------------------------------------------------

\subsection*{\khatt الاختبار 4: سؤال خارج المهمة}

اكتب:

\begin{center}

\fbox{
\parbox{11cm}{
ما هي طريقة تحضير البيتزا؟
}
}

\end{center}

لاحظ هل يحافظ المساعد على دوره.


%================================================
% النشاط السادس
%================================================

\section*{\khatt\underline{النشاط السادس: اكتشاف الأخطاء}}

سجل نتائج الاختبارات:

\begin{center}

\begin{tabular}{|
>{\centering\arraybackslash}m{2cm}|
>{\raggedleft\arraybackslash}m{5cm}|
>{\raggedleft\arraybackslash}m{5cm}|}

\hline

\textbf{الاختبار}
&
\textbf{النتيجة}
&
\textbf{المشكلة}
\\

\hline

1
&
&
\\[1cm]

\hline

2
&
&
\\[1cm]

\hline

3
&
&
\\[1cm]

\hline

4
&
&
\\[1cm]

\hline

\end{tabular}

\end{center}

اختر مشكلة واحدة تريد إصلاحها:

\[
\boxed{\hspace{15cm}}
\]


%================================================
% النشاط السابع
%================================================

\section*{\khatt\underline{النشاط السابع: تحسين التعليمات}}

نفترض أن المساعد يعطي إجابات طويلة جدًا.

نريد تعديل التعليمات.

افتح:

\begin{center}

\textenglish{Instructions}

\end{center}

وأضف:

\begin{center}

\fbox{
\parbox{12cm}{
اجعل الإجابات مناسبة للطلاب المبتدئين.

استعمل جملًا قصيرة.

استعمل نقاطًا وعناوين واضحة.

لا تقدم شرحًا طويلًا إلا إذا طلب المستخدم ذلك.
}
}

\end{center}

احفظ التعديل.


%================================================
% النشاط الثامن
%================================================

\section*{\khatt\underline{النشاط الثامن: إعادة الاختبار}}

أعد الاختبار التالي:

\begin{center}

\fbox{
\parbox{11cm}{
ما هو الذكاء الاصطناعي؟

اشرحه لطالب مبتدئ.
}
}

\end{center}

قارن الإجابة قبل التعديل وبعده.

\begin{center}

\begin{tabular}{|
>{\raggedleft\arraybackslash}m{5cm}|
>{\raggedleft\arraybackslash}m{4cm}|
>{\raggedleft\arraybackslash}m{4cm}|}

\hline

\textbf{العنصر}
&
\textbf{قبل التعديل}
&
\textbf{بعد التعديل}
\\

\hline

طول الإجابة
&
&
\\[0.8cm]

\hline

وضوح الإجابة
&
&
\\[0.8cm]

\hline

مناسبة للمبتدئ
&
&
\\[0.8cm]

\hline

تنظيم الإجابة
&
&
\\[0.8cm]

\hline

\end{tabular}

\end{center}


%================================================
% CLAUDE
%================================================

\section*{\khatt\underline{النشاط التاسع: بناء نفس المساعد في Claude}}

إذا لم تتمكن من إنشاء GPT في ChatGPT،
يمكن تنفيذ الفكرة باستعمال
\textenglish{Claude Projects}.


%------------------------------------------------

\subsection*{\khatt الخطوة 1: فتح Claude}

\begin{enumerate}

    \item افتح متصفح الإنترنت.

    \item افتح \textenglish{Claude}.

    \item سجل الدخول.

\end{enumerate}


%------------------------------------------------

\subsection*{\khatt الخطوة 2: فتح Projects}

من القائمة اختر:

\begin{center}

\textenglish{Projects}

\end{center}

ثم:

\begin{center}

\textenglish{New Project}

\end{center}


%------------------------------------------------

\subsection*{\khatt الخطوة 3: تسمية المشروع}

اكتب:

\begin{center}

\textbf{مساعد مادة الذكاء الاصطناعي}

\end{center}


%------------------------------------------------

\subsection*{\khatt الخطوة 4: إضافة تعليمات المشروع}

افتح:

\begin{center}

\textenglish{Project Instructions}

\end{center}

ثم الصق:

\begin{center}

\fbox{
\parbox{13cm}{

أنت مساعد تعليمي متخصص في مادة الذكاء الاصطناعي.

مهمتك مساعدة الطلاب على فهم الدروس بطريقة بسيطة
ومنظمة.

اشرح المفاهيم خطوة بخطوة.

استعمل أمثلة بسيطة.

اعتمد على الملفات الموجودة في معرفة المشروع
عندما يكون السؤال متعلقًا بها.

لا تخترع معلومات غير موجودة في الملفات.

إذا لم تجد المعلومة، أخبر الطالب بذلك بوضوح.

نظم الإجابات باستعمال العناوين والنقاط.

إذا كان السؤال غير واضح، اطلب توضيحًا.

}
}

\end{center}


%------------------------------------------------

\subsection*{\khatt الخطوة 5: إضافة الملف}

افتح:

\begin{center}

\textenglish{Project Knowledge}

\end{center}

ثم:

\begin{enumerate}

    \item اختر إضافة ملف.

    \item حدد نفس ملف الدرس.

    \item انتظر انتهاء رفع الملف.

    \item تأكد من ظهور الملف.

\end{enumerate}


%------------------------------------------------

\subsection*{\khatt الخطوة 6: إنشاء المحادثة}

افتح محادثة جديدة داخل المشروع.

اكتب:

\begin{center}

\fbox{
\parbox{11cm}{
اعتمد على الملف الموجود في المشروع.

ما أهم ثلاثة مفاهيم في الدرس؟

اشرحها لطالب مبتدئ.
}
}

\end{center}


%------------------------------------------------

\subsection*{\khatt الخطوة 7: اختبار المساعد}

اختبر:

\begin{enumerate}

    \item سؤالًا عاديًا.

    \item سؤالًا يعتمد على الملف.

    \item سؤالًا بمعلومات ناقصة.

    \item سؤالًا خارج المهمة.

\end{enumerate}

سجل النتيجة:

\[
\boxed{\hspace{15cm}}
\]


%================================================
% WORKFLOW
%================================================

\section*{\khatt\underline{النشاط العاشر: بناء سير عمل}}

سنقوم الآن ببناء \textenglish{Workflow}
لتحويل درس طويل إلى أسئلة تدريبية.


%------------------------------------------------

\subsection*{\khatt الخطوة 1: تحديد المدخل والمخرج}

\textbf{المدخل:}

\begin{center}

\fbox{\textbf{نص محاضرة طويل}}

\end{center}

\textbf{المخرج:}

\begin{center}

\fbox{\textbf{أسئلة تدريبية مع الإجابات}}

\end{center}

سنقسم العمل إلى ثلاث مراحل.


%------------------------------------------------

\subsection*{\khatt الخطوة 2: المرحلة الأولى}

المهمة:

\begin{center}

\textbf{استخراج الأفكار الرئيسية}

\end{center}

استخدم النص التالي:

\begin{center}

\fbox{
\parbox{12cm}{
ضع هنا نص المحاضرة التي تريد تحليلها.
}
}

\end{center}

ثم اكتب:

\begin{center}

\fbox{
\parbox{12cm}{
اقرأ النص التالي.

استخرج أهم الأفكار والمفاهيم.

اكتب كل فكرة في نقطة مستقلة.

لا تضف معلومات غير موجودة في النص.
}
}

\end{center}


%------------------------------------------------

\subsection*{\khatt الخطوة 3: حفظ النتيجة}

بعد الحصول على النتيجة:

\begin{enumerate}

    \item اقرأ النتيجة.

    \item قارنها بالنص الأصلي.

    \item صحح أي خطأ.

    \item احتفظ بالنسخة النهائية.

\end{enumerate}

\textbf{مهم:}

لا تنتقل إلى المرحلة الثانية قبل التحقق من
المرحلة الأولى.


%------------------------------------------------

\subsection*{\khatt الخطوة 4: المرحلة الثانية}

خذ نتيجة المرحلة الأولى.

ثم اكتب:

\begin{center}

\fbox{
\parbox{12cm}{
اعتمد فقط على الأفكار التالية.

اكتب ملخصًا منظمًا للدرس.

استعمل عناوين واضحة ونقاطًا قصيرة.

لا تضف معلومات جديدة.
}
}

\end{center}

احفظ الملخص.


%------------------------------------------------

\subsection*{\khatt الخطوة 5: مراجعة الملخص}

تحقق من:

\begin{itemize}

    \item هل الأفكار الرئيسية موجودة؟

    \item هل توجد معلومات إضافية غير مطلوبة؟

    \item هل توجد أخطاء؟

    \item هل الملخص مفهوم؟

\end{itemize}


%------------------------------------------------

\subsection*{\khatt الخطوة 6: المرحلة الثالثة}

خذ الملخص النهائي.

اكتب:

\begin{center}

\fbox{
\parbox{12cm}{
اعتمد فقط على الملخص التالي.

أنشئ خمسة أسئلة تدريبية للمبتدئين.

أنشئ ثلاثة أسئلة للفهم وسؤالين للتطبيق.

اكتب الإجابة الصحيحة تحت كل سؤال.

لا تستعمل معلومات غير موجودة في الملخص.
}
}

\end{center}


%------------------------------------------------

\subsection*{\khatt الخطوة 7: التحقق النهائي}

لكل سؤال تحقق من:

\begin{enumerate}

    \item هل السؤال مرتبط بالدرس؟

    \item هل الإجابة صحيحة؟

    \item هل يمكن العثور على الإجابة في المحتوى؟

    \item هل السؤال مناسب للمبتدئ؟

\end{enumerate}


%================================================
% النشاط الحادي عشر
%================================================

\section*{\khatt\underline{النشاط الحادي عشر: رسم سير العمل}}

أكمل سير العمل:

\begin{center}

\textbf{النص الأصلي}

$\downarrow$

\textbf{استخراج الأفكار}

$\downarrow$

\textbf{المراجعة}

$\downarrow$

\textbf{إنشاء الملخص}

$\downarrow$

\textbf{المراجعة}

$\downarrow$

\textbf{إنشاء الأسئلة}

$\downarrow$

\textbf{التحقق النهائي}

\end{center}

أكمل الجدول:

\begin{center}

\begin{tabular}{|
>{\centering\arraybackslash}m{1.5cm}|
>{\raggedleft\arraybackslash}m{4cm}|
>{\raggedleft\arraybackslash}m{4cm}|
>{\raggedleft\arraybackslash}m{4cm}|}

\hline

\textbf{المرحلة}
&
\textbf{المدخل}
&
\textbf{المعالجة}
&
\textbf{المخرج}
\\

\hline

1
&
النص
&
استخراج الأفكار
&
\\[1cm]

\hline

2
&
الأفكار
&
إنشاء الملخص
&
\\[1cm]

\hline

3
&
الملخص
&
إنشاء الأسئلة
&
\\[1cm]

\hline

\end{tabular}

\end{center}


%================================================
% النشاط الثاني عشر
%================================================

\section*{\khatt\underline{النشاط الثاني عشر: الاختبار النهائي}}

نفذ المهمة نفسها بطريقتين.

\textbf{الطريقة الأولى:}

استعمل Prompt واحدًا يطلب من الذكاء الاصطناعي:

\begin{center}

قراءة الدرس $\rightarrow$ تلخيصه $\rightarrow$
إنشاء الأسئلة

\end{center}

\textbf{الطريقة الثانية:}

استعمل الـ Workflow الذي بنيته:

\begin{center}

قراءة الدرس

$\downarrow$

استخراج الأفكار

$\downarrow$

مراجعة

$\downarrow$

تلخيص

$\downarrow$

مراجعة

$\downarrow$

إنشاء الأسئلة

\end{center}

قارن النتيجتين:

\begin{center}

\begin{tabular}{|
>{\raggedleft\arraybackslash}m{4cm}|
>{\centering\arraybackslash}m{4cm}|
>{\centering\arraybackslash}m{4cm}|}

\hline

\textbf{المعيار}
&
\textbf{Prompt واحد}
&
\textbf{Workflow}
\\

\hline

سهولة المراجعة
&
&
\\[0.8cm]

\hline

اكتشاف الأخطاء
&
&
\\[0.8cm]

\hline

تنظيم العمل
&
&
\\[0.8cm]

\hline

التحكم في المراحل
&
&
\\[0.8cm]

\hline

إمكانية إعادة الاستخدام
&
&
\\[0.8cm]

\hline

\end{tabular}

\end{center}


%================================================
% أسئلة الفهم
%================================================

\section*{\khatt\underline{أسئلة الفهم}}

\begin{enumerate}

    \item ما وظيفة حقل
    \textenglish{Instructions}؟

    \[
    \boxed{\hspace{15cm}}
    \]

    \item ما وظيفة
    \textenglish{Knowledge}؟

    \[
    \boxed{\hspace{15cm}}
    \]

    \item لماذا نضيف ملفات مرجعية إلى المساعد؟

    \[
    \boxed{\hspace{15cm}}
    \]

    \item لماذا يجب اختبار المساعد بعد إنشائه؟

    \[
    \boxed{\hspace{15cm}}
    \]

    \item ماذا نفعل إذا كانت إجابة المساعد غير مناسبة؟

    \[
    \boxed{\hspace{15cm}}
    \]

    \item لماذا نقسم المهمة الكبيرة إلى عدة مراحل؟

    \[
    \boxed{\hspace{15cm}}
    \]

    \item لماذا يجب مراجعة نتيجة كل مرحلة؟

    \[
    \boxed{\hspace{15cm}}
    \]

\end{enumerate}


%================================================
% المهمة النهائية
%================================================

\section*{\khatt\underline{المهمة النهائية}}

قم ببناء مساعد خاص بك.

يجب أن يحتوي المشروع على:

\begin{enumerate}

    \item اسم واضح.

    \item دور محدد.

    \item تعليمات واضحة.

    \item ملف مرجعي واحد على الأقل.

    \item ثلاثة اختبارات مختلفة.

    \item تعديل واحد على الأقل للتعليمات.

    \item إعادة اختبار بعد التعديل.

    \item Workflow من ثلاث مراحل.

\end{enumerate}

\textbf{اسم المساعد:}

\[
\boxed{\hspace{15cm}}
\]

\textbf{المهمة:}

\[
\boxed{\hspace{15cm}}
\]

\textbf{المنصة:}

\begin{center}

$\square$ \textenglish{ChatGPT}

\hspace{2cm}

$\square$ \textenglish{Claude}

\end{center}

\textbf{اسم الملف المرجعي:}

\[
\boxed{\hspace{15cm}}
\]

\textbf{عدد مراحل Workflow:}

\[
\boxed{\hspace{5cm}}
\]


%================================================
% خلاصة
%================================================

\begin{center}

\rule{0.9\textwidth}{0.5pt}

\vspace{0.2cm}

{\khatt\textbf{
المسار العملي للمخبر
}}

\vspace{0.2cm}

\textenglish{
Define
$\rightarrow$
Write Instructions
$\rightarrow$
Create
$\rightarrow$
Add Knowledge
$\rightarrow$
Test
$\rightarrow$
Debug
$\rightarrow$
Improve
$\rightarrow$
Build Workflow
$\rightarrow$
Validate
}

\vspace{0.2cm}

{\khatt\textbf{
المساعد الجيد لا يعتمد فقط على Prompt جيد،
بل يحتاج إلى تعليمات واضحة، معرفة مناسبة،
اختبارات، تحسين مستمر، ومراجعة للنتائج.
}}

\vspace{0.2cm}

\rule{0.9\textwidth}{0.5pt}

\end{center}

\end{document}